% -*-coding: utf-8 -*-

\defaultfont

\BiChapter{水下绳驱机械臂结构研制}{Some Tricks of Using this Template}
\label{Tricks}

\BiSection{引言}{Something About Templates}
\label{tricks:abouttemplates}

本章面向其水下抓取作业需求，研制一款水下绳驱机械臂。该机械臂采用绳驱动结构，将质量较大的驱动电机后置于基座，从而有效减小运动过程中对运载器的扰动。为此，本章围绕绳驱机械臂的结构设计展开研究，具体内容包括：总体方案设计、本体结构设计、驱动绳走线分析、驱动密封舱设计。

\BiSection{水下绳驱机械臂设计要求}{Something About Templates}
针对搭载于微小型运载器的对海底作业的绳驱机械臂，在设计时需要达到的以下四个设计要求：

(1) 机械臂安装于运载器的正下方，能够完全折叠；

(2) 机械臂完全伸展下探深度要大于500mm；

(3) 机械臂运动部分重量不大于3Kg；

(4) 机械臂承载能力1Kg的有效负载。

\BiSection{机械臂总体结构设计}{Something About Templates}
\BiSubsection{机械臂总体方案设计}{Necessary to Use Templates？}
\label{tricks:necessaryornot}
为满足设计要求，本文提出一种水下绳驱机械臂的初步结构方案。该方案将所有驱动部件集成于同一密封舱内，以绳索为传动介质，将舱内电机输出的力矩通过预设的绳索路径传递至各关节，从而缩小机械臂的体积与质量，降低机械臂运动时对运载器的扰动。绳驱机械臂的自由度越多，其末端精度与负载能力越低；同时，在作业过程中可通过运载器与机械臂的协同运动保证操作灵活性。综合考虑上述因素，本文选定3自由度构型，其中关节1为水平旋转关节，关节2为肘关节，关节3为腕关节。绳驱机械臂结构简图如图3.1所示。

图3.1 绳驱机械臂机构简图

为实现设计要求中机械臂完全折叠至运载器底部的功能，常规铰链关节难以满足折叠需求。为此，本文对肘部关节进行改进，采用关联双铰链滚动结构，扩展肘关节折叠空间，使大臂在非工作时段能够完全收拢于载体正下方，从而降低载体在目标搜寻过程中的水阻力。基于上述思路，本文设计的可折叠绳驱机械臂结构简图如图3.2所示。
\BiSubsection{机械臂本体结构设计}{Which Templates I should Select？}
\BiSubsection{驱动绳走线及导向设计}{Which Templates I should Select？}
\BiSubsection{驱动系统密封舱设计}{Which Templates I should Select？}
\BiSection{水下绳驱机械臂样机集成}{Something About Templates}
\BiSection{本章小结}{Something About Templates}
%\fi

%%% Local Variables: 
%%% mode: latex
%%% TeX-master: "../main"
%%% End: 
